\section{各工具评定结果的比较}
\label{sec:fatigue-compare}

将五类工具的控制工况并列，可看出工具族疲劳性能的分布规律，如表~\ref{tab:fatigue-compare} 所列。为便于比较，每种工具只列其最不利的工况。

\begin{table}[htbp]
\centering
\caption{五类工具疲劳评定结果的比较（控制工况，$n=200$ 次）}
\label{tab:fatigue-compare}
\small
\setlength{\tabcolsep}{4pt}
\begin{tabular}{@{}llcccccc@{}}
\hline
工具 & 材料档 & 控制工况 & $\sigma_{\max}$/MPa & $\sigma_{a,\mathrm{eq}}$/MPa & $N$/次 & $D$ & $n_f$/$n_s$ \\
\hline
\multirow{3}{*}{三合一} & 4145 & 拉伸 & 340.1 & 206.4 & $5.7\times10^{5}$ & $3.4\times10^{-4}$ & 2.62/2.43 \\
 & 4145H & 拉伸 & 392.3 & 246.2 & $2.3\times10^{5}$ & $8.8\times10^{-4}$ & 2.27/2.10 \\
 & 4145H & 抗扭 & 464.5 & 305.6 & $9.4\times10^{4}$ & $2.1\times10^{-3}$ & 1.92/1.78 \\
四合一 & 4330V & 拉伸 & 588.3 & 411.2 & $3.54\times10^{4}$ & $5.96\times10^{-3}$ & 1.66/1.58 \\
强磁吸附 & 4330V & 抗扭 & 597.7 & 597.7 & $6.9\times10^{3}$ & $2.9\times10^{-2}$ & 1.56/1.56 \\
过滤 & 4145H & 拉伸 & 600.6 & 413.9 & $3.7\times10^{4}$ & $5.4\times10^{-3}$ & 1.66/1.55 \\
可变径 & 4330V & 抗扭 & 627.9 & 627.9 & $9.5\times10^{3}$ & 0.021 & 1.69/1.70 \\
\hline
\end{tabular}

\vspace{2pt}
{\footnotesize 注：$D$ 为 200 趟累积损伤；$n_f$ 为疲劳安全系数，$n_s$ 为静强度安全系数；可变径工具的抗扭工况取退刀槽净截面探针值评定，五类工具的控制工况均满足判据。}
\end{table}

评定结果随载荷等级呈现明显的分层。400\,kN 级的三合一工具评定应力为 340$\sim$465\,MPa，计算寿命在 $9.4\times10^{4}\sim9.1\times10^{5}$ 次之间，寿命裕度为 470$\sim$4\,550 倍，三项指标均留有充裕余量。1\,000\,kN 级工具的评定应力集中在 527$\sim$659\,MPa，计算寿命为 $6.9\times10^{3}\sim3.7\times10^{4}$ 次，静强度安全系数为 1.54$\sim$1.62，距判据 1.5 仅 3\%$\sim$8\%。这一分布说明 1\,000\,kN 级工具的承力截面已接近充分利用：在额定载荷下各承力件的静强度安全系数被压缩到判据下限附近，若继续提高额定载荷，须同步加大截面或提高材料档。

五类工具的控制工况均满足判据，差别主要在寿命裕度。表~\ref{tab:fatigue_compare} 中可变径工具的抗扭工况裕度最紧，200 趟累积损伤 0.021、寿命裕度约 47 倍，而族内其余工具为 150$\sim$185 倍；其静强度与疲劳安全系数为 1.70 与 1.69，与其余工具（1.54$\sim$1.66）同处一档。该工况的评定应力 627.9\,MPa 取自刮刀芯轴 $\phi58$ 退刀槽净截面（1\,506.5\,mm$^2$）上的三处探针，按等效应力换算，17\,kN$\cdot$m 扭矩在该截面上相当于 $946$\,kN 的轴向拉力，与 1\,000\,kN 的额定拉力同量级，故扭转与拉伸的余量相近（1.69 与 1.70）。该结果源于退刀槽净截面小于同端螺纹根部截面这一几何特征，而非材料强度不足；距奇异点约 3\,mm 的凹角邻域仍存在 921\,MPa 的局部峰值，安全系数降至 1.16，因此加工时须保证退刀槽与台肩过渡处的圆角不小于 R1；若圆角无法保证，则应按峰值口径把额定扭矩限制在 12.9\,kN$\cdot$m 以内，或把净截面放大至 $\phi63.5$、螺纹同步放大为 Tr68$\times$4。

从取值角度看，横向比较的可靠性取决于三处取值的一致性：材料强度一律取图纸下限、评定应力一律按剔除奇异点后的探针读数、循环特征一律按载荷的实际变化取定。三者若有任一项在各工具之间取法不一，表中的数值便不能直接比较。本文采用的正是第~\ref{sec:fatigue-3in1}~节所述的统一口径，故表中各行的差异可以归因于结构、载荷与几何特征本身，而非取数方法的不同。

由上述评定可以归纳出工具族疲劳设计的三条措施。第一，螺纹根部与台阶过渡处的过渡圆角是决定寿命的首要几何因素，各工具的控制部位几乎都落在这一类特征上，加工时应保证 R3 以上的圆角，且该圆角不得在建模与网格划分中被简化。第二，配合端面与固定约束端面的峰值应力多属边界条件引起的伪应力，不应作为判据，但可作为该处对扭矩或装配预紧敏感的提示，宜加密网格复核并控制上扣扭矩。第三，成品应按 NB/T\,47013.4 进行磁粉检测，重点检查螺纹根部、退刀槽与台肩过渡处；使用中则应避免大锤敲击一类会在这些部位引入冲击应力的操作。
